\chapter{Objetivos Específicos}
\label{cap:objetivos}

Os objetivos específicos organizam-se em duas frentes complementares: a construção e a validação técnica do artefato (engenharia) e a avaliação pedagógica do treinamento com usuários.

\section{Objetivos de Engenharia}

\begin{itemize}
    \item Mapear os principais riscos associados a trabalhos em altura e ao uso inadequado de EPIs em canteiros de obras brasileiros, com base nas diretrizes da NR-18 \cite{brasil2025nr18} e NR-35 \cite{brasil2023nr35}.
    \item Projetar e implementar uma arquitetura modular orientada a eventos que desacople a lógica de negócio do motor de jogo (\textit{engine}), permitindo a autoria de cenários por dados e oferecendo afordância para múltiplos domínios de treinamento.
    \item Validar tecnicamente o sistema de forma independente do motor de jogo, por meio de suíte de testes e da execução da lógica fora do \textit{engine}, evidenciando modularidade, testabilidade e resiliência a atualizações de \gls{SDK}.
\end{itemize}

\section{Objetivos de Avaliação Pedagógica}

\begin{itemize}
    \item Avaliar o ganho de conhecimento sobre procedimentos de segurança em trabalho em altura e uso de EPIs, por meio de testes aplicados antes e após a intervenção e, condicionalmente, de um teste de retenção, comparando o grupo experimental ao grupo controle submetido a método tradicional (\textbf{H1}).
    \item Avaliar a motivação e o engajamento dos participantes nas dimensões do modelo ARCS, por meio do RIMMS, comparando os dois grupos (\textbf{H2}).
    \item Avaliar a usabilidade e a sensação de presença do sistema junto ao grupo experimental, por meio do \gls{SUS} e do \gls{IPQ} (\textbf{H3}).
\end{itemize}
